\documentclass{article}

\usepackage{amsmath,amssymb}
\usepackage{xurl}
\usepackage[hidelinks]{hyperref}
\setlength{\emergencystretch}{2em}

\input{command}

\title{The Missing Order Hypothesis in Wolf's Unitary Comparison}
\author{}
\date{2026-08-24}

\begin{document}
\maketitle
\gapnote{false-source}{resolved}

\section{The assertion}

Let $A,B\in\MN{d}$ be Hermitian, and write
$\lambda_1^\downarrow(A),\ldots,\lambda_d^\downarrow(A)$ for the eigenvalues
of $A$ in decreasing order. The paragraph preceding Wolf's Proposition~5.3
assumes $A\leq B$ and applies Weyl's monotonicity theorem in the form
\begin{align}
    \lambda_j^\downarrow(B)
        &\geq \lambda_j^\downarrow(A)
            +\lambda_d^\downarrow(B-A),
        \qquad 1\leq j\leq d.
    \label{eq:wolf_ch5_unitary_comparison_missing_order_weyl_bound}
\end{align}
This is Equation~(5.56) of Ref.~\cite{Wolf2012Quantum}. Since $B-A\geq0$, its
smallest eigenvalue is nonnegative. Hence
\begin{align}
    \lambda_j^\downarrow(A)
        &\leq \lambda_j^\downarrow(B),
        \qquad 1\leq j\leq d.
    \label{eq:wolf_ch5_unitary_comparison_missing_order_ordered_eigenvalue_comparison}
\end{align}

Proposition~5.3 then drops the hypothesis $A\leq B$. It claims that arbitrary
Hermitian $A,B\in\MN{d}$ admit a unitary $U$, independent of $f$, such that
\begin{align}
    f(A)
        &\leq U f(B)U^\dagger
    \label{eq:wolf_ch5_unitary_comparison_missing_order_printed_unitary_comparison}
\end{align}
for every non-decreasing function $f:I\to\mathbb R$, where $I$ is an interval
containing both spectra. This is Equation~(5.57) of
Ref.~\cite{Wolf2012Quantum}.\footnote{Source:
    \cite[Proposition~5.3 and Equations~(5.56)--(5.57)]
    {Wolf2012Quantum}; local source
    \path{Notes/WolfNoteTexSource/ch05_schwarz_inequalities.tex},
    lines~1119--1135. Tracked in \ghissue{27}.}

\section{The point at issue}

The printed proposition is false. In dimension one, set
\begin{align}
    A&=[1],
    &
    B&=[0],
    &
    I&=[0,1],
    &
    f(x)&=x.
    \label{eq:wolf_ch5_unitary_comparison_missing_order_scalar_counterexample}
\end{align}
Every one-dimensional unitary is a scalar phase, so
\begin{align}
    U f(B)U^\dagger
        &=0,
    \label{eq:wolf_ch5_unitary_comparison_missing_order_scalar_conjugation}\\
    f(A)\leq U f(B)U^\dagger
        &\Longleftrightarrow 1\leq0.
    \label{eq:wolf_ch5_unitary_comparison_missing_order_scalar_contradiction}
\end{align}
Thus the hypotheses printed in Proposition~5.3 of
Ref.~\cite{Wolf2012Quantum} do not imply its conclusion.

The source proof also requires the omitted order assumption. Equation~(5.56)
of Ref.~\cite{Wolf2012Quantum} assumes $A\leq B$, and this assumption yields
the ordered eigenvalue comparison
(\ref{eq:wolf_ch5_unitary_comparison_missing_order_ordered_eigenvalue_comparison}).
Without it, the proof supplies no comparison between the two ordered spectra.

\section{The corrected statement and source proof}

The source argument proves the following corrected proposition:
\begin{align}
    A\leq B
        \quad\Longrightarrow\quad
        \exists U\in\MN{d}\ \Bigl[
        &U^\dagger U=UU^\dagger=\Id
    \label{eq:wolf_ch5_unitary_comparison_missing_order_corrected_unitary_quantifier}\\
        &{}\land
        \forall I\subseteq\mathbb R\
        \forall f:I\to\mathbb R,
        \nonumber\\
        &\bigl(        I\text{ is an interval}
            \ \land\
            \operatorname{spec}(A)\cup\operatorname{spec}(B)\subseteq I
            \ \land\
            f\text{ is non-decreasing on }I
        \bigr)
        \nonumber\\
        &\hspace{4em}\Longrightarrow
            f(A)\leq U f(B)U^\dagger
        \Bigr].
    \label{eq:wolf_ch5_unitary_comparison_missing_order_corrected_statement}
\end{align}
The unitary is chosen before the interval and the function, so the same $U$
works for every admissible $I$ and $f$.

Choose unitaries $V_A$ and $V_B$ whose columns are eigenvectors arranged in
decreasing eigenvalue order. Then
\begin{align}
    A
        &=V_A\operatorname{diag}\bigl(        \lambda_1^\downarrow(A),\ldots,
            \lambda_d^\downarrow(A)
        \bigr)V_A^\dagger,
    \label{eq:wolf_ch5_unitary_comparison_missing_order_A_diagonalization}\\
    B
        &=V_B\operatorname{diag}\bigl(        \lambda_1^\downarrow(B),\ldots,
            \lambda_d^\downarrow(B)
        \bigr)V_B^\dagger.
    \label{eq:wolf_ch5_unitary_comparison_missing_order_B_diagonalization}
\end{align}
The order hypothesis and Equation~(5.56) of
Ref.~\cite{Wolf2012Quantum} give
(\ref{eq:wolf_ch5_unitary_comparison_missing_order_ordered_eigenvalue_comparison}).
Monotonicity of $f$ therefore implies
\begin{align}
    f(\lambda_j^\downarrow(A))
        &\leq f(\lambda_j^\downarrow(B)),
        \qquad 1\leq j\leq d.
    \label{eq:wolf_ch5_unitary_comparison_missing_order_functional_eigenvalue_comparison}
\end{align}
Set
\begin{align}
    U
        &=V_A V_B^\dagger.
    \label{eq:wolf_ch5_unitary_comparison_missing_order_eigenbasis_unitary}
\end{align}
Functional calculus in the two eigenbases gives
\begin{align}
    f(A)
        &=V_A\operatorname{diag}\bigl(        f(\lambda_1^\downarrow(A)),\ldots,
            f(\lambda_d^\downarrow(A))
        \bigr)V_A^\dagger,
    \label{eq:wolf_ch5_unitary_comparison_missing_order_functional_calculus_A}\\
    U f(B)U^\dagger
        &=V_A\operatorname{diag}\bigl(        f(\lambda_1^\downarrow(B)),\ldots,
            f(\lambda_d^\downarrow(B))
        \bigr)V_A^\dagger.
    \label{eq:wolf_ch5_unitary_comparison_missing_order_functional_calculus_B}
\end{align}
The diagonal comparison
(\ref{eq:wolf_ch5_unitary_comparison_missing_order_functional_eigenvalue_comparison}),
transported by the unitary congruence $X\mapsto V_AXV_A^\dagger$, proves
(\ref{eq:wolf_ch5_unitary_comparison_missing_order_corrected_statement}).

\section{Formalization}

The cross-matrix Weyl comparison is formalized by
\leanid{LinearMap.IsSymmetric.eigenvalues\_le\_of\_sub\_isPositive}, with
matrix forms \leanid{Matrix.IsHermitian.eigenvalues₀\_mono} and
\leanid{Matrix.IsHermitian.eigenvalues\_mono}.\footnote{Formal module:
    \doclink{QICLean/Analysis/WeylMonotonicity}.}
Its finite-dimensional variational proof intersects the span of the first
$j+1$ eigenvectors of $A$ with the span of the last $d-j$ eigenvectors of
$B$. Their dimensions sum to $d+1$, so they contain a common nonzero vector
$x$. The ordered spectral expansions and positivity of $B-A$ yield
\begin{align}
    \lambda_j^\downarrow(A)\lVert x\rVert^2
        &\leq \operatorname{Re}\braket{Ax}{x},
    \label{eq:wolf_ch5_unitary_comparison_missing_order_variational_lower_bound}\\
    \operatorname{Re}\braket{Ax}{x}
        &\leq \operatorname{Re}\braket{Bx}{x},
    \label{eq:wolf_ch5_unitary_comparison_missing_order_positive_difference_bound}\\
    \operatorname{Re}\braket{Bx}{x}
        &\leq \lambda_j^\downarrow(B)\lVert x\rVert^2.
    \label{eq:wolf_ch5_unitary_comparison_missing_order_variational_upper_bound}
\end{align}
Since $x\neq0$, these inequalities prove the required cross-matrix
eigenvalue comparison.

The corrected proposition is formalized by
\leanid{Matrix.IsHermitian.exists\_unitary\_cfc\_le\_cfc\_of\_le}.\footnote{Formal
    module: \doclink{QICLean/Channel/Schwarz/UnitaryComparison}.}
The proof chooses the unitary from the two decreasingly ordered eigenbases and
uses the formalized Weyl comparison for every interval and every
non-decreasing function.

\section{Conclusion}

Wolf's Proposition~5.3 is false as printed because it omits $A\leq B$. The
scalar example
(\ref{eq:wolf_ch5_unitary_comparison_missing_order_scalar_counterexample})
satisfies the printed hypotheses and contradicts the conclusion. Adding
$A\leq B$ makes the proposition follow from Equation~(5.56) of
Ref.~\cite{Wolf2012Quantum}: compare the ordered spectra, apply scalar
monotonicity of $f$, and align the eigenbases with one unitary. Both the
cross-matrix Weyl comparison and this corrected unitary comparison are
formalized without assuming a pointwise eigenvalue comparison.

\bibliographystyle{alpha}
\bibliography{references}

\end{document}
